\section{Search-o1 的结果、边界与 Search-R1}

\subsection{随着文档增加，性能仍能扩展}

\videofigure{figures/search-o1-scaling.jpg}{Search-o1 在多个科学领域随检索文档数量增加持续提升，而标准 RAG 很快饱和。}{01:00:34--01:01:28}

\videofigure{figures/gpqa-results.jpg}{Search-o1 在 GPQA 扩展集上超过直接推理和标准 RAG，并缩小与人类专家的差距。}{01:01:28--01:03:18}

\videofigure{figures/retrieved-doc-scaling.jpg}{从 top-1 扩展到 top-10 文档时，Search-o1 能利用更多证据；普通 RAG 更容易被噪声拖累。}{01:03:18--01:04:38}

\videofigure{figures/multihop-qa.jpg}{在 HotpotQA、2Wiki、MuSiQue 与 Bamboogle 等多跳 QA 上，迭代搜索优于单跳检索。}{01:04:38--01:06:12}

这些结果支持一个系统设计原则：检索规模要与证据提炼能力共同扩展。只增加文档数量而不提高 query refinement 和 compression，无法获得稳定收益。

\subsection{不确定性减少不等于事实完全可靠}

\videofigure{figures/search-o1-insights.jpg}{论文从不确定词、文档规模与多轮查询三个角度分析 Search-o1 的收益来源。}{01:06:12--01:06:46}

\videofigure{figures/reasoning-quality.jpg}{加入搜索后，推理中的模糊措辞减少，证据支持的陈述增加。}{01:06:46--01:08:02}

\begin{warningbox}{“说得更确定”不是正确性的充分条件}
检索器可能返回错误网页，压缩器可能误读证据，模型也可能把引用与结论错误绑定。仍需来源质量、claim-level citation 与最终答案验证。
\end{warningbox}

\subsection{Search-o1 与 Search-R1}

Search-o1 主要通过 prompting 和显式 agent loop 教模型何时搜索；Search-R1 则尝试用 reinforcement learning 把搜索策略学进 policy。二者差异可概括为：

\begin{table}[H]
\centering
\small
\begin{tabular}{>{\raggedright\arraybackslash}p{3cm}>{\raggedright\arraybackslash}p{4.7cm}>{\raggedright\arraybackslash}p{4.7cm}}
\toprule
\textbf{维度} & \textbf{Search-o1} & \textbf{Search-R1} \\
\midrule
控制方式 & prompt、特殊 token 与手写循环 & RL 学习何时搜索和如何使用结果 \\
优点 & 易调试、可替换检索组件、无需训练 & 可减少手工规则，策略可能更自适应 \\
风险 & prompt 脆弱、调用策略未必最优 & reward 设计复杂，可能学会投机搜索 \\
适用阶段 & 快速构建与分析系统 & 有可靠 reward 和足够在线交互预算 \\
\bottomrule
\end{tabular}
\end{table}

\begin{importantbox}{搜索行为本身也需要 reward}
若只奖励最终答案，agent 可能过度搜索、完全不搜或利用 benchmark 漏洞。理想目标还应考虑答案质量、证据覆盖、查询次数、延迟和文档成本。
\end{importantbox}

\subsection{本章小结}

Search-o1 证明多轮、按需、可压缩的检索能显著改善知识密集推理；Search-R1 进一步把调用策略视为可学习行为。两者都要求验证检索证据，而不是把 retrieval 当作天然真相。

